\section{Discussion}
\label{sec:discussion}

\subsection{Principal findings}

\paragraph{The architecture transfers.}
INTERVenE-Enc, tuned and evaluated on MIMIC-IV ICU-level data
\citep{oded2027intervene}, retains its ordering against strong raw-stream
baselines on a very different cohort: multicenter, general ward, older,
longer stays, and shorter observation window in relative terms (K = 4 of
14 days here versus K = 2 of 14 days on MIMIC). No hyperparameters were
tuned to this cohort. Support-weighted AUPRC lifts from $0.472$ for
ss-STraTS to $0.527$ for INTERVenE-Enc on KBTA, with non-overlapping
bootstrap intervals---a similar relative gap to the MIMIC-IV finding
($0.631 \to 0.672$).

\paragraph{The knowledge base does more here than on MIMIC-IV.}
The single most informative contrast in the paper is the $\sigma$-bin
arm against the KB arm ($5 \to 6$): identical Transformer, identical
data, identical labels, replacing only the abstraction layer with a
seven-band $\sigma$-discretization. On MIMIC-IV that substitution was
within noise \citep{oded2027intervene}; on Kineret it costs $0.073$
AUPRC ($0.454 \to 0.527$) and $0.050$ AUROC, and drops the arm to the
logistic-regression floor. The interpretation is that a knowledge base
authored for the diabetes-inpatient network it is deployed on carries
predictive information that a knowledge-free discretization does not
recover---even when the discretization is fed to the same
Transformer.

Two related observations follow. First, this is where clinical domain
knowledge and empirical modelling meet materially: the same authors
constructed the abstraction rules \citep{hochberg2023inpatient} that
give this arm its lift and are unavailable to $\sigma$-bins. Second,
knowledge-base contribution to accuracy is cohort-dependent, and
studies of the value of clinical knowledge in EHR modelling that
report a single dataset (MIMIC or otherwise) should be read
accordingly.

\paragraph{Measured guideline adherence does not shift aggregate
metrics.}
The QA contribution is bounded by $\le 0.008$ on any aggregate metric
across three architectures (Figure~\ref{fig:qa-delta}). The signal is
used---sparse-transcoder attribution places compliance-pattern tokens
among the top drivers for seven of eight outcomes---but is largely
redundant with the interval stream a KB-consuming model is already
reading, because both are deterministic derivations of the same
underlying abstraction engine. The concurrent adherence study
\citep{oded2027kineret-jbi} identified above-median adherence as
associated with lower mortality (OR $=0.557$ after propensity
weighting); those population-level associations do not need to
translate to per-admission predictive lift for a model already reading
the underlying abstractions.

This has a design consequence. If a downstream task cares about the
score \emph{as an operational input}---e.g., a dashboard alerting on
low compliance---the QA layer stands on its own on its aggregation
target. If the task is complication prediction from a model that
already consumes the abstractions, no additional AUPRC accrues from
supplying the score. Which of the two the deployment cares about
determines whether the QA layer is a predictor input, a monitoring
readout, or both.

\subsection{The design that let both questions be answered}

Two decisions distinguish this study from an ordinary architecture
comparison.

\paragraph{One canonical event set.}
Complications appear in two source tables: the raw table
(\texttt{mediator\_input.csv}, containing the measurements the
knowledge base reads) and the Mediator output (containing its verdicts).
These disagree on which occurrences happened and when, both because
measurement-derived rules were not authored twice and because timestamps
drift by minutes across derivations. We rebuild the raw side's own
complication rows in Python by applying the same threshold rules the
Mediator's XML applies, then intersect the two views at $1$-minute
tolerance and keep only occurrences the two files date identically. The
result is one canonical event set that supplies labels, is re-injected
into every arm's input stream, and is asserted by a unit test to be
identical across arms. Without this step, on this cohort, the raw file
carried zero lab-threshold complications and the Mediator carried
$68{,}525$; cross-file alignment of the pre-computed complications
found a support drift of $3{,}096$ vs $3{,}098$ / $2{,}849$ vs $2{,}854$
that would otherwise have made \emph{which events did each arm see} a
confound rather than a constant.

\paragraph{Knowledge containment.}
The Mediator emits \texttt{*\_EVENT} concepts that are not prediction
targets---\texttt{AKI\_EVENT}, \texttt{ELECTROLYTE\_DERANGEMENT\_EVENT},
\texttt{MYOCARDIAL\_INJURY\_EVENT}. Each is a knowledge-base verdict
reached by reading raw measurements. Handing them to LogReg or
ss-STraTS would give those arms clinical judgements they could not have
computed themselves, and rungs $5 \to 6$ would then measure a leak
rather than the abstraction lift. The two streams differ by design:
non-target \texttt{*\_EVENT} concepts are stripped from the raw stream
and kept in the KB stream; target events go to everyone as observed
history. This is the one place a wiring mistake would silently
invalidate every cross-arm number in the paper, so it is asserted by a
unit test in both directions.

\subsection{Where the knowledge base helps most}

The KB arm's largest per-outcome lifts are on the two frequent glycemic
outcomes---hyperglycemia $0.578$ over ss-STraTS $0.486$ ($+0.092$
AUPRC); severe hyperglycemia $0.513$ over $0.472$ ($+0.041$). These are
threshold events by definition: hyperglycemia is a glucose level above
a defined cut, and the knowledge base collapses hundreds of lab codes
into a single \emph{glucose-state} channel with the threshold applied.
The model reads the clinical definition of its label directly, which is
close to what an interpretable clinical scoring rule would do.

Death is more interesting: no single KBTA threshold defines it, and the
lift ($0.645$ vs $0.559$ for ss-STraTS) suggests the model composes
evidence across several abstracted states. This is consistent with the
sparse-transcoder attribution reported on MIMIC-IV, which found
hypoalbuminemia (a chronic-mortality marker), aminotransferase
elevation (liver stress), and bicarbonate shifts (acid--base
decompensation) as the top death drivers. The same clinical logic---
composition across abstracted trends rather than reading a single
threshold---applies to LoS ($44.4$\,h vs $\sim 48$\,h on MIMIC-IV,
$45.9$\,h vs $46.4$\,h here).

Below roughly $200$ positives, the arms stop separating: infection,
MI/cardiovascular, and severe hypoglycemia sit near their prevalence
floors regardless of arm, and the single largest swing in the per-
outcome table---the $\sigma$-bin arm at $0.295$ on MI/cardiovascular
against $0.564$ for logistic regression---rests on $69$ events and
should not be read as an ordering. This is a study-power finding rather
than a modelling one: the cohort under-samples rare complications in
absolute terms, and running a larger extract would sharpen those rows
without needing any change to the pipeline.

\subsection{Limitations}
\label{sec:discussion-limits}

\paragraph{One cohort, single-seed training.}
Every arm is trained once, at the reference-repository hyperparameters,
without a per-cohort tune. Multi-seed variance was not affordable
inside the regulated environment; the reported bootstrap intervals
capture sampling variability for the fitted model, not the variance of
the training procedure. This is the same limitation reported by the
methodological MIMIC-IV study \citep{oded2027intervene}.

\paragraph{External validity is limited.}
This is the second cohort INTERVenE has been evaluated on, and both are
diabetic. Whether the architecture-transferred and knowledge-lift-larger
findings extend to non-diabetic complication prediction, or to
non-Israeli hospital networks, is not settled here---though the design
of the ladder is portable and the release of the code
(Section~\ref{sec:reproducibility}) makes the experiment repeatable
elsewhere.

\paragraph{The knowledge base itself is not evaluated here.}
The knowledge base used to produce the KB stream is the diabetes-
inpatient ontology of \citep{hochberg2023inpatient,oded2027kineret-jbi};
this paper takes it as given. Its predictive contribution is what rungs
$5 \to 6$ measure, but the ontology's coverage, calibration and edge
cases are the subject of the concurrent adherence study, not this one.

\paragraph{Class imbalance at the tail of the target list.}
Four of the eight targets have fewer than $200$ positives on the test
set. Their intervals overlap across most arms and their per-outcome
comparisons should not be over-read. This is not a modelling artefact
but a data-scale one; a larger extract would resolve it.

\paragraph{Model calibration is not reported.}
Reliability curves and expected calibration error were computed on the
MIMIC-IV cohort and reported per outcome in
\citep{oded2027intervene}. We omit the corresponding analysis here to
stay inside the AIIM figure budget, but the run-outputs contain the
per-prediction probabilities needed to reproduce them
(Section~\ref{sec:reproducibility}); this is~\TODO{a candidate for the
supplementary material.}

\paragraph{No temporal or per-hospital held-out.}
The split is patient-level and random. The concurrent adherence study
demonstrated a small but measurable shift in adherence patterns after
the 2023 guideline was issued, and a temporal-holdout evaluation
(train pre-guideline, test post-guideline, or vice versa) would test
whether prediction quality tracks that shift. Similarly, a per-hospital
leave-one-out would test whether the model generalizes across sites
within the network. Both are~\TODO{listed as open reviewer questions in
Section~\ref{sec:tasks-open} of the draft appendix.}
